% Maintained chapter; figures integrated from reviewed e9e5ddea assets.
\section{问题重述与分析}
\label{sec:problem_analysis}

\subsection{问题背景}

山区洪涝灾害发生后，道路中断、桥梁受损以及局部通信基础设施失效，会显著降低传统地面救援的可达性与时效性。无人机不依赖既有道路网络，能够快速跨越复杂地形，适合承担应急物资运输与临时通信保障任务。然而，在山区环境中，无人机任务不仅受载荷、航程和地形净空约束，还同时受到货箱不可拆分、实体无人机与电池数量有限、任务截止时间、通信遮挡以及资源独立配置等条件影响。

本题给定一个运输基地、15个救援服务区、80个不可拆货箱、三类运输无人机以及相应能源与通信资源，要求依次研究单架次安全运输、多任务联合调度、运输--通信协同和独立任务组资源配置。四个问题使用同一组地形、货物与装备数据，后续问题在前一问题基础上逐步增加新的现实约束，因此需要建立统一的底层物理模型，并保证各问之间的数据口径和任务继承关系一致。

\subsection{问题重述}

\subsubsection{问题一：安全运输能力与货箱组批}

在只考虑单点往返运输的条件下，需要结合服务区位置、DEM地形、无人机载荷--航程关系和返航安全余量，确定三类无人机到各服务区执行任务时的最大安全载荷。在此基础上，对各服务区的不可拆货箱进行组批并选择运输机型，使全部货箱均能安全送达，同时尽量减少运输架次，并兼顾能耗和累计作业时间。

\subsubsection{问题二：有限无人机与共享电池条件下的联合调度}

在问题一单架次物理模型的基础上，进一步允许一次任务访问多个服务区，并考虑实体无人机数量、共享电池库存、返航后的电池恢复以及货箱送达时限。需要同时确定货箱组合、访问顺序、机型、实体无人机、电池和任务开始时刻，使运输计划真实可执行，并在满足配送时限的前提下尽可能缩短全部运输任务的完成时间。

\subsubsection{问题三：连续通信约束下的运输--中继协同}

山区地形可能使运输无人机在飞行过程中与地面站之间发生通信遮挡。问题三要求在问题二运输任务结构的基础上，为直连失效区间配置中继无人机，并综合考虑中继位置与高度、到达时间、服务窗口、中继实体和能源组件等条件，使运输全过程满足通信要求，同时优化运输与中继共同构成的系统完成时间和能耗。

\subsubsection{问题四：独立任务组的资源配置}

在问题三正式运输--通信方案确定后，将15个服务区划分为若干相互独立的任务组。不同任务组之间不能共享运输无人机、电池、中继无人机和中继能源组件，需要判断哪些服务区由于既定运输或中继任务而必须保持在同一组，并计算各组独立执行时每类资源的最小配置量以及相对现有库存的资源缺口。

\subsection{问题分析}

\subsubsection{问题一的分析}

问题一首先需要解决“额定载荷”与“山区实际安全载荷”之间的差异。对于同一机型，不同服务区的水平距离、航段最高地形和爬升高度不同；同时无人机航程又随载荷非线性变化，因此最大安全载荷应由完整往返能耗约束反求，而不能直接使用额定载荷。

得到各“服务区--机型”的安全载荷后，问题转化为带质量、体积和能量约束的不可拆货箱组批。由于每个服务区的货箱数量有限，可先枚举可行批型，再建立精确覆盖模型。本文采用原始箱级MILP、对称压缩MILP与动态规划进行交叉求解，以最少架次为首要目标，并在固定架次数后继续比较总能耗和累计作业时间；同时设置启发式基线，用于衡量精确组批带来的改进。

\subsubsection{问题二的分析}

问题二的难点由“一个任务是否可飞”转变为“全部任务能否在有限资源下按时执行”。一项运输任务不仅占用实体无人机，还会占用与机型兼容的电池；无人机返航后即可进入下一任务，而电池还需继续充电恢复，因此两类资源具有不同的占用区间。此外，多服务区路线会改变逐段载荷、送达时刻和任务持续时间，使货箱组合、路径和排程相互耦合。

因此，本文先由统一物理模型生成满足质量、体积、能量和路线要求的运输任务，再对实体无人机、电池与开始时刻进行联合排程。目标采用分层处理：首先保证加权迟到最小，在零迟到条件下进一步压缩系统最后返航时间。考虑到原问题组合规模较大，仅给出一个高质量可行解不足以说明解的质量，故求解同时设置两条路线：一条持续改进可行上界，另一条通过任务列松弛、完整定价和分支覆盖构造全局有效下界，从而给出当前方案的确定性最优性区间。

\subsubsection{问题三的分析}

问题三新增的核心约束不是静态“覆盖到某一区域”，而是运输无人机沿连续轨迹飞行时，在每个需要通信的时刻都必须存在可用链路。因此需要先根据真实三维运输轨迹识别直连失效区间，再判断候选中继在完整服务区间内能否同时满足运输机--中继和中继--地面站链路条件。

当前正式模型继承问题二的23个运输任务结构，在此基础上重新安排运输开始时刻，并联合优化中继位置、高度、服务窗口、中继实体和能源组件。中继往返转场采用当前正式方案的兼容悬停端点高度规则，普通运输航段的地形净空规则保持不变。求解采用“候选生成--完整几何重建--联合排程--固定离散结构连续时间精修--独立验证”的闭环：离散候选只负责产生较优结构，最终接受的方案均返回真实轨迹，对完整通信区间与临界端点重新核验。由于当前搜索没有穷尽连续三维选址空间，本问给出经过完整验证的当前最优已知联合方案，不扩大为原连续问题的全局最优结论。

\subsubsection{问题四的分析}

问题四不重新设计问题三的运输和中继计划，而是研究既定救援预案被划分为独立执行单元后的资源代价。若同一运输架次访问多个服务区，或同一逻辑中继任务实际保障多个服务区，则这些服务区不能被拆到不同任务组，否则就必须改变已经冻结的任务。因此应首先依据运输访问关系和实际通信保障关系构造服务区关联图，将15个服务区压缩为不可拆任务块。

在固定绝对时刻后，每一类资源的最小配置问题可转化为区间最大重叠问题。对每个候选分区，分别计算各组内A/B/C型运输机、电池、中继无人机和能源组件的最大并发占用数，再将各组需求相加并与库存比较。当前正式日程只形成3个不可拆块，因此两组和三组情形的全部严格合法分区数量分别为3和1，可以直接完整枚举，不需要启发式搜索。

\subsection{总体建模思路}

四问虽然优化目标不同，但共享同一套空间、地形、载荷、时间和能耗计算基础。本文首先在第3章建立统一运输任务物理模型，将任意候选任务转换为可直接用于后续优化的距离、巡航高度、逐段载荷、任务时长、能耗和箱级送达偏移；随后按照“能力确定--资源调度--通信协同--独立配置”的顺序逐层增加约束。

四问之间传递的任务属性与冻结边界见图\ref{fig:decision_inheritance}。

\CommonDecisionFigure

问题一给出单架次能力边界和精确组批结果；问题二将可行运输任务放入真实无人机--电池时间轴，并利用上下界证据评价工期解质量；问题三进一步把通信保障嵌入运输时间轴，在完整区间上验证中继可执行性；问题四冻结问题三正式日程，研究组织边界改变后异构资源共享能力的损失。这样既避免四问重复建立物理模型，又保证后续每一层优化都能追溯到同一套底层数据和计算口径。

\FloatBarrier
